\documentclass[../../main.tex]{subfiles}
\graphicspath{{../../image/}{../../figure/}}
\begin{document}

\section{伪激波总强度的界}\label{section:ch14-伪激波总强度的界}

验证 $U_\delta$ 是\eqref{align:ch14-system}的 $\delta$-近似解的最后任务是建立要求\eqref{align:ch14-pseudoshock-strength}。生成阶概念已在\cref{section:ch14-全局波型}中引入。高生成阶波在大量碰撞之后产生，因而预期其强度很小。确实，论证的第一步是证明：所有波（特别是所有伪激波）中生成阶充分高者的总强度任意小。为此，按\cref{section:ch14-总变差的界}的分析加以精细化，把波按其生成阶分别筛选并监控。

我们现已知道碰撞总数有界，因而所有波的生成阶落在有限范围 $0\leqslant\mu\leqslant\nu$。但应注意 $\nu$ 的大小归根结底依赖于 $\delta$，并预期当 $\delta\downarrow0$ 时无界增长。对 $\mu=0,1,\cdots,\nu$ 与 $t\in[0,\infty)\setminus I$，令 $L_\mu(t)$ 表示穿过 $t$-时间线且生成阶 $\geqslant\mu$ 的所有波的强度之和；$Q_\mu(t)$ 表示穿过 $t$-时间线且生成阶 $\mu_1,\mu_2$ 满足 $\max\{\mu_1,\mu_2\}\geqslant\mu$ 的所有接近波对强度之积之和。特别地，$L_0(t)=L(t)$ 且 $Q_0(t)=Q(t)$。最后，把 $I$ 中生成阶 $\mu$ 的波与生成阶 $\leqslant\mu$ 的波碰撞的时刻集合记为 $I_\mu$。

生成阶 $\leqslant\mu-2$ 的波之间的碰撞不影响生成阶 $\geqslant\mu$ 的波，因而
\begin{align}
\Delta L_\mu(t)=0,\qquad t\in I_0\cup\cdots\cup I_{\mu-2}.
\label{align:ch14-deltaLmu0}
\end{align}
$L_\mu(\cdot)$ 在 $t\in I$ 的任一变化必由两个波碰撞诱发，其中至少一个生成阶 $\geqslant\mu-1$。这些碰撞波（强度记为 $|\alpha|$ 与 $|\beta|$）对 $Q_{\mu-1}(t-)$ 贡献 $|\alpha||\beta|$，而对 $Q_{\mu-1}(t+)$ 无贡献。如\cref{section:ch14-总变差的界}，$Q_{\mu-1}(\cdot)$ 的相应下降可用于抵消 $L_\mu(\cdot)$ 的潜在增量，后者以 $\kappa|\alpha||\beta|$ 为界：
\begin{align}
\Delta L_\mu(t)+2\kappa\Delta Q_{\mu-1}(t)\leqslant0,\qquad t\in I_{\mu-1}\cup\cdots\cup I_\nu.
\label{align:ch14-deltaLmu}
\end{align}

由类似论证可验证不等式
\begin{align}
\Delta Q_\mu(t)+2\kappa\Delta Q(t)L_\mu(t-)\leqslant0,\qquad t\in I_0\cup\cdots\cup I_{\mu-2},
\label{align:ch14-deltaQmu1}
\end{align}
\begin{align}
\Delta Q_\mu(t)+2\kappa\Delta Q_{\mu-1}(t)L(t-)\leqslant0,\qquad t\in I_{\mu-1},
\label{align:ch14-deltaQmu2}
\end{align}
\begin{align}
\Delta Q_\mu(t)\leqslant0,\qquad t\in I_\mu\cup\cdots\cup I_\nu,
\label{align:ch14-deltaQmu3}
\end{align}
它们支配 $Q_\mu(\cdot)$ 跨各阶碰撞的变化。

下面用上标 $+$ 或 $-$ 表示``正''或``负''部分：$w^+=\max\{w,0\}$，$w^-=\max\{-w,0\}$。目的是监控量
\begin{align}
\hat L_\mu=\sup_t L_\mu(t),\qquad \hat Q_\mu=\sum_{t\in I}[\Delta Q_\mu(t)]^+,
\label{align:ch14-hats}
\end{align}
对 $\mu=1,\cdots,\nu$，并证明
\begin{align}
\hat L_\mu\leqslant2^{-\mu}c_3a,\qquad \hat Q_\mu\leqslant2^{-\mu+3}c_3^2a^2,
\label{align:ch14-hat-bounds}
\end{align}
其中 $a$ 是 $TV_{(-\infty,\infty)}U_0(\cdot)$ 的界。

由\eqref{align:ch14-deltaLmu0}、\eqref{align:ch14-deltaLmu}与``初始条件''$L_\mu(0+)=0$，$\mu=1,\cdots,\nu$，可得
\begin{align}
\hat L_\mu\leqslant2\kappa\sum_{t\in I}[\Delta Q_{\mu-1}(t)]^-,\qquad \mu=1,\cdots,\nu.
\label{align:ch14-hatL}
\end{align}

接着关注\eqref{align:ch14-deltaQmu1}、\eqref{align:ch14-deltaQmu2}与\eqref{align:ch14-deltaQmu3}，以 $Q_\mu(0+)=0$ 为``初始条件''。回忆\eqref{align:ch14-Q-bound}、\eqref{align:ch14-L-bound}并利用
\begin{align}
\sum_{t\in I}[\Delta Q(t)]^-=Q(0+)-Q(\infty)\leqslant\tfrac12 L(0+)^2,
\label{align:ch14-Q-negative}
\end{align}
我们推得
\begin{align}
\hat Q_\mu\leqslant\kappa L(0+)^2\hat L_\mu+4\kappa L(0+)\sum_{t\in I}[\Delta Q_{\mu-1}(t)]^-,\qquad \mu=1,\cdots,\nu.
\label{align:ch14-hatQ}
\end{align}

结合\eqref{align:ch14-hatL}与\eqref{align:ch14-hatQ}。设初值总变差充分小，使 $10\kappa L(0+)\leqslant1$，则
\begin{align}
\hat Q_\mu\leqslant\frac12\sum_{t\in I}[\Delta Q_{\mu-1}(t)]^-,\qquad \mu=1,\cdots,\nu.
\label{align:ch14-hatQ-half}
\end{align}
特别地，对 $\mu=1$ 并由\eqref{align:ch14-Q-negative}，我们推得 $\hat Q_1\leqslant\frac14 L(0+)^2$。

最后注意，对 $\mu=1,\cdots,\nu$，因 $Q_\mu(0+)=0$，
\begin{align}
\sum_{t\in I}[\Delta Q_\mu(t)]^-=\sum_{t\in I}[\Delta Q_\mu(t)]^+-Q_\mu(\infty)\leqslant\hat Q_\mu.
\label{align:ch14-Qmu-negative}
\end{align}
因而\eqref{align:ch14-hatQ-half}给出 $\hat Q_\mu\leqslant\frac12\hat Q_{\mu-1}$，$\mu=2,\cdots,\nu$，进而 $\hat Q_\mu\leqslant2^{-\mu-1}L(0+)^2$。连同\eqref{align:ch14-Q-negative}与\eqref{align:ch14-hatQ}给出估计 $\hat L_\mu\leqslant2^{-\mu-2}L(0+)$。由此建立了\eqref{align:ch14-hat-bounds}。

现显然可取充分大的 $\mu_0$，使生成阶 $\geqslant\mu_0$ 的所有波总强度（以 $\hat L_{\mu_0}$ 为上界）不超过 $\frac12\delta$。

为估计生成阶 $<\mu_0$ 的伪激波总强度，第一步是估计其数目。对 $\mu=0,\cdots,\nu$，令 $K_\mu$ 表示生成阶 $\leqslant\mu$ 的波数。$K_\mu$ 的粗上界可由如下论证导出。通过两个 Riemann 解算器之一求解跳跃间断所产生的出射波数以数 $b/\delta$ 为界。因而 $K_0\leqslant\frac{b}{\delta}N$，其中 $N$ 是 $U_{0\delta}$ 的跳跃点数。由于任两波在其生存期内至多碰撞一次，可能产生生成阶 $\mu$ 波的碰撞数以 $\frac12 K_{\mu-1}^2$ 为界。因而
\begin{align}
K_\mu\leqslant K_{\mu-1}+\frac{b}{2\delta}K_{\mu-1}^2\leqslant\frac{b}{\delta}K_{\mu-1}^2,
\label{align:ch14-Kmu}
\end{align}
由此易得
\begin{align}
K_\mu\leqslant\left(\frac{b}{\delta}\right)^{2^{\mu+1}}N^{2^\mu}.
\label{align:ch14-Kmu-bound}
\end{align}

接着估计单个伪激波的强度。任一伪激波由两个前沿（强度 $|\alpha|$ 与 $|\beta|$ 满足 $|\alpha||\beta|\leqslant\sigma$）碰撞产生，因而经简化 Riemann 解算器求解，如\cref{figure:ch14-不同族前沿碰撞}与\cref{figure:ch14-同族前沿碰撞}所示。由相应相互作用估计\eqref{align:ch14-pseudoshock-gen}可知，任一伪激波在诞生时强度不超过 $c_4\sigma$。由\eqref{align:ch14-pseudoshock-hit}，伪激波与强度 $|\beta|$ 的前沿（如\cref{figure:ch14-伪激波与前沿碰撞}所示）碰撞至多使其强度增大因子 $1+\kappa|\beta|$。因而伪激波的强度最终至多增大因子 $\prod(1+\kappa|\gamma|)$，其中乘积遍历该伪激波在生存期内碰撞的所有前沿。由于伪激波非共振，估计\eqref{align:ch14-LC-bound}适用并蕴含 $\sum|\gamma|\leqslant2L(0+)$。设 $4\kappa L(0+)\leqslant1$，我们得到任一伪激波在任一时刻的强度不超过 $3c_4\sigma$。

现显然通过采用\eqref{align:ch14-Kmu-bound}所给 $K_{\mu_0-1}$ 的上界，并选取充分小的 $\sigma$，可保证生成阶 $<\mu_0$ 的伪激波总强度以 $\frac12\delta$ 为界。结合先前关于生成阶 $\geqslant\mu_0$ 伪激波总强度的估计，这建立了\eqref{align:ch14-pseudoshock-strength}。

\end{document}
